\section{Discussion}

\subsection{Dynamic Feedback Control}

ODCR treats explanation degeneration as a dynamic causal-control problem. D4C's
feedback view explains why domain textual shortcuts can become
self-reinforcing: \(W^{(t)}\) influences explanations, and explanations can
feed back into \(W^{(t+1)}\). ODCR does not discard this diagnosis. It refines
the feedback loop so that stable content evidence, dynamic style, and
counterfactual reliability are controlled separately.

\subsection{Role of EASD and HSS}

EASD and HSS are the parts of ODCR that keep the paper from becoming centered
on reliability routing alone. EASD anchors content/style decomposition in evidence, separating what
should be explained from how it is expressed. HSS makes style feedback
hierarchical, with domain-global style and instance-local residuals. Together,
they define where dynamic feedback should operate: generated explanations may
reshape style, but they should not overwrite content evidence.

\subsection{Role of RCR/UCI and CCV/FCA}

RCR/UCI controls counterfactual influence after decomposition. It filters
counterfactual candidates, scores reliability, and uses uncertainty to reduce
trust in unstable shifts. CCV/FCA then turns the controlled evidence into
target-domain explanations. CCV conditions generation on \(U\), \(I\), \(C\),
\(S\), and \(R_{cf}\); FCA keeps explanation content aligned with
content-grounded rating evidence. This closes the ODCR loop from dynamic
feedback diagnosis to controlled faithful generation.

\subsection{Reading the Current Evidence}

The current result tables provide a compact main-text comparison over four
directed transfer settings. The Amazon blocks illustrate related-domain
content/style transfer, while the TripAdvisor--Yelp blocks place more stress on
style shift and reliability control. These results support the manuscript as
bounded current evidence. Additional ablations, human assessment, multi-seed
evaluation, and broader baseline adaptation can be integrated as later evidence
without changing the core method story.
